\section{تدقيق عامل القياس والكثافات المحلية}

هذا القسم جزء من التدقيق الرياضي لما بعد التأليف. وهو يصحح ثلاثة مواضع كانت في المسودة الأولى بحاجة إلى ضبط: اعتماد التكامل المفرد على $N$، وشرط إيجابية السلسلة المفردة، وتصنيف مساهمة الأقواس الكبرى.

\subsection{فصل عامل القياس عن الثابت الأرخميدي}

نضع
\[
w_k(\gamma)=\int_0^1 e(\gamma u^k)\,du.
\]
وبتغيير المتغير $t=Pu$ في تعريف
\[
v_k(\beta;P)=\int_0^P e(\beta t^k)\,dt
\]
نحصل على
\[
v_k(\beta;P)=P\,w_k(\beta P^k).
\]
وبما أن $P^k=N$، فإن تغيير المتغير $\gamma=\beta P^k$ يعطي الهوية الدقيقة
\[
\mathfrak J_{s,k}(N)
=
\int_{-\infty}^{\infty}v_k(\beta;P)^s e(-N\beta)\,d\beta
=
P^{s-k}\mathcal J_{s,k},
\]
حيث
\[
\mathcal J_{s,k}
=
\int_{-\infty}^{\infty}w_k(\gamma)^s e(-\gamma)\,d\gamma.
\]
إذن العامل الذي يحمل النمو مع $N$ هو
\[
P^{s-k}=N^{s/k-1},
\]
أما $\mathcal J_{s,k}$ فهو الثابت الأرخميدي بعد التطبيع. وعندما يكون $s>k$ يكون التكامل متقاربًا تقاربًا مطلقًا من الحد القياسي
\[
w_k(\gamma)\ll_k \min\{1,|\gamma|^{-1/k}\}.
\]
وفي نطاق وارينغ الكلاسيكي تكون قيمته الموجبة
\[
\mathcal J_{s,k}
=
\frac{\Gamma(1+1/k)^s}{\Gamma(s/k)}.
\]

\begin{proposition}[هوية قياس التكامل المفرد]
\label{prop:chapter17-singular-integral-scaling}
إذا كان $P^k=N$ و$\mathfrak J_{s,k}(N)$ معرفًا كما سبق، فإن
\[
\mathfrak J_{s,k}(N)=N^{s/k-1}\mathcal J_{s,k}.
\]
وعند $s>k$ يكون $\mathcal J_{s,k}$ متقاربًا تقاربًا مطلقًا وموجبًا.
\resultid{ANT-PROP-17-02}
\provedhere
\end{proposition}

\begin{proof}
الهوية تتبع مباشرة من التغييرين $t=Pu$ و$\gamma=\beta P^k$. أما التقارب المطلق فينتج من $|w_k(\gamma)|\ll_k\min\{1,|\gamma|^{-1/k}\}$، لأن $|w_k(\gamma)|^s$ قابل للتكامل عند اللانهاية إذا كان $s>k$. والإيجابية تتبع من التفسير الحجمي القياسي للتكامل بوصفه كثافة أرخميدية لمساحة الحلول الموجبة للمعادلة $u_1^k+\cdots+u_s^k=1$.
\end{proof}

\subsection{ما الذي يكفي لإيجابية السلسلة المفردة؟}

لا يكفي أن نقول فقط إنه «لا يوجد عائق توافق ظاهر». الشرط القياسي الأقوى هو وجود حل غير منفرد في $\mathbb Z_p$ لكل أولي $p$؛ أي حل $p$-أدي لا تنعدم عنده جميع المشتقات الجزئية المناسبة. تحت هذا الشرط، وفي نطاق عدد متغيرات يضمن تقارب حاصل الضرب المحلي، تكون الكثافة المحلية $\sigma_p(N)$ موجبة، وتكتب السلسلة المفردة على الصورة
\[
\mathfrak S_{s,k}(N)=\prod_p \sigma_p(N)>0.
\]
أما إثبات التقارب المطلق لهذا الحاصل وإيجابيته بانتظام في النطاق الكلاسيكي فيبقى مدخلًا مقتبسًا من نظرية وارينغ، لا نتيجة داخلية جديدة في هذا الفصل.

\begin{remark}[التصنيف الحاكم]
تقارب السلسلة المفردة وإيجابيتها في المبرهنة الكلاسيكية مصنفان
\[
\texttt{CITED / LOCAL-DENSITY INPUT}.
\]
المثبت هنا هو فقط الفصل المنطقي بين: وجود حلول محلية، وعدم الانفراد، وتقارب حاصل الضرب، وإيجابية الحد الرئيس.
\end{remark}

\subsection{التصنيف النهائي للأقواس الكبرى}

الصيغة
\[
\int_{\mathfrak M}f_k(\alpha;P)^s e(-N\alpha)\,d\alpha
=
\mathfrak S_{s,k}(N)\mathfrak J_{s,k}(N)
+o(P^{s-k})
\]
ليست مثبتة كاملة داخل هذا الفصل، لأن توحيد خطأ التقريب المحلي، وقطع السلسلة والتكامل، وتمديدهما إلى اللانهاية تعتمد على تقديرات كلاسيكية مقتبسة. لذلك يكون تصنيف `ANT-PROP-17-01` النهائي في هذه النسخة:
\[
\texttt{CITED / EXPLAINED},
\]
لا `PROVED-HERE` ولا `AUTHORED-DRAFT` مستقلًا.

وبهذا تصبح الصيغة الكلاسيكية لوارينغ أكثر وضوحًا:
\[
r_{s,k}(N)
=
\mathfrak S_{s,k}(N)\mathcal J_{s,k}N^{s/k-1}
+o\!\left(N^{s/k-1}\right),
\]
في نطاق $s$ الكلاسيكي الكافي، ومع الفروض المحلية اللازمة.
